% =========================================================
% Chapter 6 -- ARCH and GARCH models
% Author : Aymen Ben Brik (aymen.benbrik@esprit.tn)
% =========================================================

\chapter{ARCH and GARCH models}
\label{ch:arch-garch}

\section*{Introduction}
\addcontentsline{toc}{section}{Introduction}

The ARMA family of Chapter~\ref{ch:arma} captures the dynamics of
the \emph{conditional mean}, but assumes a constant variance. On
financial returns, on macroeconomic surprises and on many physical
measurements, this assumption is violated: large absolute deviations
cluster in time, and the variance itself is predictable.
This chapter introduces the \emph{autoregressive conditional
heteroskedasticity} (ARCH) family that models the conditional
variance directly, generalises it to GARCH$(p, q)$, derives the
stationarity conditions, presents the maximum-likelihood estimator,
and applies the model to volatility forecasting and Value-at-Risk.

% =========================================================
\section{Stylised facts of financial returns}
\label{sec:stylised}

\begin{rembox}[Empirical regularities]
For daily log-returns $r_t = \log(P_t / P_{t-1})$ on broad equity
indices, currencies and commodities:
\begin{enumerate}
  \item $\E[r_t] \approx 0$ (small drift).
  \item ACF of $r_t$: nearly flat -- returns are approximately
        serially uncorrelated.
  \item ACF of $r_t^2$ (or $|r_t|$): \emph{strongly} positive at
        many lags -- volatility is autocorrelated.
  \item Marginal kurtosis $\widehat\kappa > 3$: heavier tails than
        Gaussian.
  \item Leverage effect: negative shocks raise future volatility
        more than positive shocks of the same size.
\end{enumerate}
\end{rembox}

\begin{rembox}[Conditional vs.\ unconditional]
Let $\mathcal F_{t-1} = \sigma(r_{t-1}, r_{t-2}, \dots)$ be the
information set up to $t - 1$.
\begin{itemize}
  \item \emph{Unconditional variance}:
        $\sigma^2 = \Var(r_t)$, the long-run, average risk.
  \item \emph{Conditional variance}:
        $\sigma_t^2 = \Var(r_t \mid \mathcal F_{t-1})$, the
        short-run risk knowing the past.
\end{itemize}
ARCH and GARCH models \emph{specify $\sigma_t^2$} while keeping
$\sigma^2$ finite -- the process is unconditionally homoskedastic
but conditionally heteroskedastic.
\end{rembox}

\begin{defbox}[Mean--variance decomposition]
\label{def:mean-var}
A second-order martingale-difference innovation
$\varepsilon_t = r_t - \mu_t$ is written
\[
  \varepsilon_t = \sigma_t\, z_t,
  \qquad
  z_t \iid \mathcal D(0, 1),
  \qquad
  \E[\varepsilon_t \mid \mathcal F_{t-1}] = 0,
  \quad
  \Var(\varepsilon_t \mid \mathcal F_{t-1}) = \sigma_t^2.
\]
$\mathcal D$ is the \emph{innovation distribution}, typically
$\mathcal N(0, 1)$ or Student-$t$ with low degrees of freedom.
\end{defbox}

% =========================================================
\section{ARCH$(q)$ model}
\label{sec:arch}

\begin{defbox}[ARCH$(q)$, Engle 1982]
\label{def:arch}
The \emph{autoregressive conditional heteroskedasticity} model of
order $q$ specifies
\[
  \varepsilon_t = \sigma_t z_t,
  \qquad
  \sigma_t^2 = \alpha_0
              + \sum_{i = 1}^{q} \alpha_i\, \varepsilon_{t-i}^2,
\]
with $z_t \iid \mathcal N(0, 1)$, $\alpha_0 > 0$, and
$\alpha_i \geq 0$ for $i = 1, \dots, q$.
\end{defbox}

\begin{propbox}[Properties of ARCH$(q)$]
\label{prop:arch}
Assume the covariance-stationarity condition
$\sum_{i = 1}^{q} \alpha_i < 1$. Then:
\begin{enumerate}
  \item $\E[\varepsilon_t] = 0$ and $(\varepsilon_t)$ is a white
        noise (uncorrelated in time, but not independent).
  \item $\E[\varepsilon_t^2] = \alpha_0 / (1 - \sum \alpha_i)
        =: \sigma^2$.
  \item $(\varepsilon_t^2)$ admits an AR$(q)$ representation
        $\varepsilon_t^2 = \alpha_0
                          + \sum \alpha_i \varepsilon_{t-i}^2 + \nu_t$,
        with $\nu_t = \varepsilon_t^2 - \sigma_t^2$ a martingale
        difference.
  \item Excess kurtosis: under Gaussian innovations
        ARCH$(1)$ has kurtosis
        $3\,(1 - \alpha_1^2) / (1 - 3 \alpha_1^2) > 3$ when
        $\alpha_1^2 < 1/3$.
\end{enumerate}
\end{propbox}

\begin{rembox}[Why squares?]
Squaring past residuals makes the response \emph{sign-blind}: a
$-3 \sigma$ surprise and a $+3 \sigma$ surprise both raise
$\sigma_{t+1}^2$. To allow asymmetric responses (leverage), one
moves to GJR-GARCH or EGARCH (Section~\ref{sec:extensions}).
\end{rembox}

% =========================================================
\section{GARCH$(p, q)$ model}
\label{sec:garch}

In practice, ARCH$(q)$ requires many lags $q$ to capture the
slow decay of squared-return autocorrelations, leading to long,
poorly identified parameter vectors. Bollerslev (1986) generalised
the model by including past conditional variances themselves.

\begin{defbox}[GARCH$(p, q)$, Bollerslev 1986]
\label{def:garch}
\[
  \sigma_t^2 = \alpha_0
              + \sum_{i = 1}^{q} \alpha_i\, \varepsilon_{t-i}^2
              + \sum_{j = 1}^{p} \beta_j\, \sigma_{t-j}^2,
\]
with $\alpha_0 > 0$, $\alpha_i, \beta_j \geq 0$.
\end{defbox}

\begin{propbox}[Stationarity of GARCH$(p, q)$]
\label{prop:garch-stat}
$(\varepsilon_t)$ admits a covariance-stationary solution with
finite unconditional variance if and only if
\[
  \sum_{i = 1}^{q} \alpha_i + \sum_{j = 1}^{p} \beta_j \;<\; 1.
\]
The unconditional variance is then
\[
  \sigma^2 = \E[\varepsilon_t^2]
          = \frac{\alpha_0}{1 - \sum_i \alpha_i - \sum_j \beta_j}.
\]
\end{propbox}

\begin{exbox}[GARCH$(1, 1)$: the workhorse]
\label{ex:garch11}
\[
  \sigma_t^2 = \alpha_0
             + \alpha_1\, \varepsilon_{t-1}^2
             + \beta_1\, \sigma_{t-1}^2.
\]
On daily equity returns, typical estimates are $\alpha_1 \approx
0.05$--$0.10$, $\beta_1 \approx 0.85$--$0.92$,
$\alpha_1 + \beta_1 \approx 0.95$--$0.99$. Volatility is
\emph{very persistent} but still mean-reverting.
\end{exbox}

\begin{rembox}[IGARCH and persistence]
The boundary case $\alpha_1 + \beta_1 = 1$ defines an
\emph{IGARCH}: shocks have a permanent effect on the conditional
variance. The unconditional variance is then infinite (the model
is strictly stationary but not covariance-stationary). Estimated
$\alpha_1 + \beta_1$ very close to $1$ is a sign that the assumed
GARCH$(1, 1)$ may be missing a structural break.
\end{rembox}

% =========================================================
\section{Maximum-likelihood estimation}
\label{sec:garch-mle}

\begin{defbox}[Conditional Gaussian likelihood]
\label{def:garch-mle}
Assuming $z_t \iid \mathcal N(0, 1)$ and $r_t = \mu + \varepsilon_t$,
the conditional log-likelihood is
\[
  \log \mathcal L(\bm\theta)
  = -\tfrac{T}{2}\log(2 \pi)
    -\tfrac{1}{2} \sum_{t = 1}^{T}
        \Bigl(
            \log \sigma_t^2(\bm\theta)
            + \frac{(r_t - \mu)^2}{\sigma_t^2(\bm\theta)}
        \Bigr),
\]
with $\sigma_t^2$ computed recursively from the GARCH equation.
$\bm\theta = (\mu, \alpha_0, \alpha_1, \dots, \alpha_q,
\beta_1, \dots, \beta_p)$.
\end{defbox}

\begin{thbox}[Asymptotic properties of the QMLE]
Under regularity conditions (positivity, stationarity, bounded
fourth moment), the Gaussian quasi-MLE
$\widehat{\bm\theta}_{\text{QMLE}}$ is consistent and
asymptotically normal,
\[
  \sqrt{T}\,(\widehat{\bm\theta} - \bm\theta_0)
  \;\xrightarrow{d}\; \mathcal N(0, \Sigma),
\]
\emph{even if} the true innovation distribution is non-Gaussian
(e.g.\ Student-$t$). The asymptotic variance $\Sigma$ takes a
sandwich form $\Sigma = J^{-1} I J^{-1}$.
\end{thbox}

\begin{rembox}[Practical numerics]
\begin{itemize}
  \item Returns are typically rescaled (multiplied by $100$ to work
        in percentage units) for numerical stability.
  \item Initial $\sigma_1^2$ is set to the sample variance, or to
        the unconditional value $\alpha_0 / (1 - \alpha_1 -
        \beta_1)$ at the current parameters.
  \item In Python, the \texttt{arch} package
        (\texttt{arch\_model}) does the rescaling, parameter
        constraints and likelihood maximisation automatically.
\end{itemize}
\end{rembox}

% =========================================================
\section{Diagnostics}
\label{sec:garch-diag}

\subsection{Engle's LM test for ARCH effects}

\begin{defbox}[Engle's LM test]
\label{def:engle-lm}
Apply on the residual $\widehat\varepsilon_t$ of a mean model
(constant or ARMA). Run the auxiliary regression
\[
  \widehat\varepsilon_t^2
  = \gamma_0 + \gamma_1 \widehat\varepsilon_{t-1}^2
            + \dots + \gamma_q \widehat\varepsilon_{t-q}^2 + u_t,
\]
get its $R^2$. Test
\[
  H_0: \gamma_1 = \dots = \gamma_q = 0,
  \qquad
  LM = T \cdot R^2 \;\xrightarrow{d}\; \chi^2_q.
\]
Reject $H_0$ at level $\alpha$ if $LM > \chi^2_{q, 1 - \alpha}$.
\end{defbox}

\begin{rembox}[How to use it]
Run Engle's LM \emph{before} fitting GARCH (to confirm there is an
ARCH effect to model) and \emph{after} (on the standardised
residuals $\widehat z_t = \widehat\varepsilon_t / \widehat\sigma_t$,
to confirm the model has captured it).
\end{rembox}

\subsection{Standardised residual diagnostics}

\begin{rembox}
A well-fitted GARCH produces standardised residuals that look like
white noise both in level and in square:
\begin{itemize}
  \item ACF of $\widehat z_t$ inside Bartlett bands;
  \item ACF of $\widehat z_t^2$ inside Bartlett bands;
  \item Engle LM on $\widehat z_t^2$ does not reject;
  \item marginal distribution close to $\mathcal D(0, 1)$ (QQ-plot,
        Kolmogorov--Smirnov).
\end{itemize}
Persistent failures in $\widehat z_t^2$ suggest a higher-order GARCH
or a Student-$t$ innovation; failures only in $\widehat z_t^-$
suggest a leverage extension (GJR, EGARCH).
\end{rembox}

% =========================================================
\section{Volatility forecasting and Value-at-Risk}
\label{sec:garch-forecast}

\begin{propbox}[GARCH$(1, 1)$ multi-step forecast]
\label{prop:garch-forecast}
Let $\widehat\sigma_T^2$ and $\widehat\varepsilon_T$ be the values
obtained at the end of the sample. Then
\[
  \widehat\sigma_{T+1}^2
  = \alpha_0 + \alpha_1 \widehat\varepsilon_T^2
             + \beta_1 \widehat\sigma_T^2,
\]
and for $h \geq 2$,
\[
  \widehat\sigma_{T+h}^2
  = \sigma^2
  + (\alpha_1 + \beta_1)^{h-1}
        \bigl(\widehat\sigma_{T+1}^2 - \sigma^2\bigr),
\]
where $\sigma^2 = \alpha_0 / (1 - \alpha_1 - \beta_1)$ is the
unconditional variance. The forecast \emph{mean-reverts} to
$\sigma^2$ at rate $\alpha_1 + \beta_1$.
\end{propbox}

\begin{defbox}[Conditional Value-at-Risk]
\label{def:var}
At horizon $1$ and confidence $1 - \alpha$ (typically $\alpha = 1\%$
or $5\%$), assuming Gaussian innovations,
\[
  \mathrm{VaR}_{T+1}(\alpha)
  = -\bigl(\widehat\mu + z_\alpha\, \widehat\sigma_{T+1}\bigr),
\]
with $z_\alpha = \Phi^{-1}(\alpha)$ ($z_{0.01} = -2.326$,
$z_{0.05} = -1.645$). The interpretation: with probability
$1 - \alpha$, the loss does not exceed $\mathrm{VaR}_{T+1}(\alpha)$.
\end{defbox}

\begin{rembox}[Static vs dynamic VaR]
The classical static VaR uses the \emph{sample} standard deviation
$\widehat\sigma$. The GARCH-VaR uses
$\widehat\sigma_{T+1}$. In calm regimes, $\widehat\sigma_{T+1} <
\widehat\sigma$ and the dynamic VaR is tighter; in turbulent
regimes (just after a market shock), $\widehat\sigma_{T+1} >
\widehat\sigma$ and the dynamic VaR is wider, accurately reflecting
the elevated short-term risk.
\end{rembox}

% =========================================================
\section{Extensions (overview)}
\label{sec:extensions}

\begin{rembox}[EGARCH (Nelson 1991)]
The exponential GARCH models $\log \sigma_t^2$, removing the
positivity constraints on the parameters and naturally allowing
asymmetric responses:
\[
  \log \sigma_t^2
  = \omega + \alpha\, \bigl(|z_{t-1}| - \E |z_{t-1}|\bigr)
          + \gamma\, z_{t-1}
          + \beta\, \log \sigma_{t-1}^2.
\]
$\gamma < 0$ is the typical sign of the leverage effect on equity
returns.
\end{rembox}

\begin{rembox}[GJR-GARCH (Glosten--Jagannathan--Runkle 1993)]
Adds a threshold term that activates only on negative shocks:
\[
  \sigma_t^2
  = \alpha_0
  + (\alpha_1 + \gamma\, \mathbb 1_{\varepsilon_{t-1} < 0})
       \varepsilon_{t-1}^2
  + \beta_1\, \sigma_{t-1}^2.
\]
$\gamma > 0$ implies that bad news raises future volatility more
than good news of the same magnitude.
\end{rembox}

% =========================================================
\section*{Summary}
\addcontentsline{toc}{section}{Summary}

\noindent\textbf{Key takeaways.}
\begin{itemize}
  \item Financial returns are uncorrelated but their squares are
        autocorrelated -- volatility comes in clusters.
  \item ARCH$(q)$: $\sigma_t^2$ is a linear function of past
        squared shocks. GARCH$(p, q)$ adds memory through past
        conditional variances.
  \item GARCH$(1, 1)$ is stationary iff $\alpha_1 + \beta_1 < 1$,
        with unconditional variance
        $\alpha_0 / (1 - \alpha_1 - \beta_1)$.
  \item Engle's LM test detects ARCH effects in residuals;
        standardised residuals validate the fit.
  \item Conditional volatility forecasts mean-revert to the
        unconditional variance at rate $\alpha_1 + \beta_1$;
        Value-at-Risk uses these forecasts directly.
  \item Extensions: EGARCH and GJR-GARCH for asymmetric (leverage)
        effects.
\end{itemize}
